\documentclass{beamer}
\usepackage{amsmath}
\usepackage{gvv}

\title{Question 12.14}
\author{AI25BTECH11040 - Vivaan Parashar}
\date{\today}

\begin{document}

\frame{\titlepage}

\begin{frame}
    \frametitle{Question: }
    The eigenvalues of the matrix, $\vec{A}^{-1}$, where $\vec{A} = \myvec{2 & 1\\ 0 & 3}$ are
    \begin{multicols}{2}
        \begin{enumerate}[label=\alph*)]
            \item 1 and 1/2
            \item 1 and 1/3
            \item 2 and 3
            \item 1/2 and 1/3
        \end{enumerate}
    \end{multicols}
\end{frame}

\begin{frame}
    \frametitle{Solution: }
    The eigenvalues of a matrix $\vec{A}$ are given by the roots of the characteristic polynomial, which is defined as:
    \[\text{det}(\vec{A} - \lambda \vec{I}) = 0\]
    where $\lambda$ represents the eigenvalues and $\vec{I}$ is the identity matrix of the same size as $\vec{A}$ ($\vec{I}_2$).

    To find the eigenvalues of $\vec{A}^{-1}$, we can use the property that the eigenvalues of the inverse of a matrix are the reciprocals of the eigenvalues of the original matrix.
    This can be proven using the fact
    \begin{align*}
        \vec{A}\vec{v} = \lambda \vec{v} & \implies \vec{A}^{-1}\vec{A}\vec{v} = \vec{A}^{-1}(\lambda \vec{v})   \\
                                         & \implies \vec{v} = \lambda \vec{A}^{-1}\vec{v}                        \\
                                         & \implies \vec{A}^{-1}\vec{v} = \left(\frac{1}{\lambda}\right) \vec{v}
    \end{align*}
\end{frame}
\begin{frame}
    Therefore, we first need to find the eigenvalues of $\vec{A}$.
    The characteristic polynomial of $\vec{A}$ is given by:
    \[\text{det}(\vec{A} - \lambda \vec{I}) = \text{det}\left|\myvec{2 - \lambda & 1\\ 0 & 3 - \lambda}\right| = (2 - \lambda)(3 - \lambda) - 0 = (2 - \lambda)(3 - \lambda)\]
    Setting this equal to zero gives us the eigenvalues of $\vec{A}$:
    \[(2 - \lambda)(3 - \lambda) = 0\]
    Thus, the eigenvalues of $\vec{A}$ are $\lambda_1 = 2$ and $\lambda_2 = 3$.

    Now, the eigenvalues of $\vec{A}^{-1}$ are the reciprocals of these eigenvalues:
    \[\mu_1 = \frac{1}{\lambda_1} = \frac{1}{2}, \quad \mu_2 = \frac{1}{\lambda_2} = \frac{1}{3}\]

    Therefore, the eigenvalues of $\vec{A}^{-1}$ are $\frac{1}{2}$ and $\frac{1}{3}$.
    The correct answer is option (d): 1/2 and 1/3.
\end{frame}

\end{document}
